\section{Compiled point queries from an interval-orbit certificate}
\label{sec:compiled-orbit-index}

The maintained interval kernel implements the polynomial, binary-encoded
orbit algorithm of Agol--Hass--Thurston~\cite{aht}. Its current \texttt{same\_orbit}
interface performs two orbit counts for a pair of points: count the original
relation, then adjoin their identification and count again.  Reusing only the
original count still leaves one new reduction schedule for each distinct
query.  We give an additive interface that prepares one checked trace and
answers subsequent point queries by arithmetic transport.  This is a
data-structure refinement of the existing orbit kernel, rather than a new
orbit-counting theorem.

Compressed transport histories are established mathematical tools.
For example, Erickson--Nayyeri retain a tracing history as a straight-line
program for compressed curves~\cite{erickson-nayyeri}. The contribution
here is the explicit refinement of the maintained source-bound AHT
certificate into a reusable minimum-representative interface, together
with the compact minimum-set theorem below. We do not assert priority
for the general idea of retaining a compressed trace.

An input has $k$ partial interval isometries on $[0,N)$, with integer endpoints
and either translation or reflection on each domain.  Its components are the
equivalence classes generated by the isometries.  A complete source-bound AHT
trace contains identity deletion, reflection trimming, periodic merger,
transmission, static contraction and suffix truncation.  We first replay it
with the maintained independent verifier.  No maximality assumption is used
for the correctness of a supplied trace; maximality belongs to the producer's
polynomial progress theorem.

\subsection{Every truncation moves a point to an earlier representative}

The first four operations preserve the relation on the same point universe
and require no change of a queried point.  Suppose a suffix truncation reduces
the current universe $[0,n)$ to $[0,c)$.  For a positive carrier of period
$p>0$, its valid truncation conditions imply $c-p\ge0$, and the point map is
\[
 f(x)=\begin{cases}
 x,&x<c,\\
 c-p+((x-(c-p))\bmod p),&x\ge c.
 \end{cases}
\]
For a reflection with left domain endpoint $a$ it is
\[
 f(x)=\begin{cases}x,&x<c,\\a+n-1-x,&x\ge c.\end{cases}
\]
The latter carrier is already trimmed, so its domain lies strictly to the
left of its range.  In either case a removed point maps to a surviving point
in its old orbit, and $f(x)<x$ on the removed suffix.  Division and remainder
evaluate an arbitrarily long translation chain in one binary operation.

For a static contraction let its disjoint inclusive gaps be
$[a_i,b_i]$, in increasing order, and set
\[
 P_i=\sum_{0\le j<i}(b_j-a_j+1),\qquad P_0=0.
\]
The gaps are indexed starting at zero, so $P_i$ is the total length of
preceding gaps. Every point in a gap is currently a
singleton orbit and is emitted.  A surviving point becomes
\[
 x\longmapsto x-\sum_{b_i<x}(b_i-a_i+1).
\]
Binary search among the gap endpoints evaluates both gap membership and
this rank shift.  Its inverse on survivors inserts the preceding gap
lengths.  The collapsed breakpoints are $a_i-P_i$; a breakpoint not greater
than the survivor coordinate contributes its gap length to the inverse.
Suffix truncation needs no inverse map when recovering the original
coordinate of a surviving representative.

\begin{theorem}[Canonical minimum representatives from a checked trace]
For every complete valid AHT trace there is a compiled index $m$ such that
\[
 m(x)=\min\{y\in[0,N):y\sim x\}.
\]
It is obtained without listing points or components.  In particular,
$x\sim y$ if and only if $m(x)=m(y)$, and minimum representatives agree
between any two valid traces for the same source relation.
\end{theorem}
\begin{proof}
Follow $x$ through the displayed rank shifts and folds until it reaches
a static gap.  The local replay theorem says that surviving equivalence
classes and emitted singleton representatives together describe exactly
the previous relation; therefore two original points receive the same
terminal emission precisely when they were equivalent.

For minimality, associate each current point with its original source
coordinate.  The current source-coordinate map is increasing: a contraction
deletes points and closes gaps, while a truncation retains a prefix.
Inductively, the source coordinate associated with each current point is
the least original point transported to it.  Contraction preserves this
property.  Under truncation, every removed current point moves strictly to
the left, so all original points transported with it are larger than the
source coordinate already associated with the surviving target.  Thus the
property is preserved by folds as well.  Once a singleton is emitted, no
other live point remains in its equivalence class, and its source coordinate
is the minimum of the entire original orbit.  Reverse all preceding static
rank shifts to recover that source coordinate.  This proves the formula
and hence its independence from the chosen trace.
\end{proof}

The index also numbers emitted singleton points by a binary ordinal in
$[0,C)$, where $C$ is the orbit count, and supports selecting the minimum
representative of a supplied ordinal.  Store the cumulative number emitted
before each contraction and before each gap.  Two binary searches find its
emission coordinate; reversing earlier contractions returns its source
minimum.  This order is deliberately not claimed to be canonical.  For the
identity relation on three points with an explicit identity generator at
zero, one valid trace deletes that generator before contracting, producing
order $(0,1,2)$.  Another first contracts the static points one and two,
then deletes the generator, producing order $(1,2,0)$.  Their minimum maps
are nevertheless identical.

\subsection{The entire canonical minimum set has few ordinary intervals}

There is a stronger consequence of the one-directional truncations.
Although the orbit count can be enormous, all original minimum
representatives admit a small description as ordinary intervals; no
arithmetic-progression encoding is needed for this particular set.

\begin{theorem}[Compact canonical minimum set]
If a valid complete trace has $g$ static-gap records, its set of original
orbit minima is a union of at most $g$ disjoint integer intervals.  This
canonical union can be constructed in polynomial time from the checked
trace.  After storing prefix lengths, selection by increasing minimum
and ranking a supplied minimum require $O(\log(g+2))$ comparisons and
binary additions.  Their order is independent of the trace.
\end{theorem}
\begin{proof}
Let $M$ be the original minimum representatives emitted so far, and let
$L$ be the original source coordinates of current live points.  We maintain
the invariant
\[
 L=[0,b)\setminus M
\]
for some source cutoff $b$; $M$ is allowed to have members beyond $b$.
Initially $b=N$ and $M$ is empty.  A suffix truncation retains an initial
segment of the ordered live set, so it only decreases $b$.  A static
contraction moves some live source points from $L$ into $M$ and leaves
the same equation valid with unchanged $b$.

Consider one current static gap, whose endpoint source coordinates are
$u\le v$.  In the source interval $[u,v]$, the currently live points are
exactly those represented by this gap.  Every other point in that source
interval already belongs to $M$, by the invariant.  Thus emitting the gap
changes the minimum set by the exact update
\[
 M\longmapsto M\cup[u,v].
\]
Previously truncated nonminimum points cannot lie between $u$ and $v$:
they lie beyond the cutoff of the current live prefix.  Apply this update
once per recorded gap.  The final minimum set is therefore the union of
at most $g$ source intervals.  Lift each gap's endpoints by reversing only
the preceding static contractions, sort their source hulls, and merge
overlap and adjacency.  A prefix sum of the resulting interval lengths
gives both stated rank/selection queries by binary search.  The first
theorem identifies this set intrinsically as all orbit minima, establishing
its trace independence.
\end{proof}

The interval bound is sharp.  Take $r$ disjoint consecutive pairs of
blocks of width $w$, with the translation identifying the two blocks in
each pair.  There are $rw$ components, whose minima form exactly
\[
 \bigcup_{i=0}^{r-1}[2iw,(2i+1)w).
\]
A right-to-left valid trace has one static gap per block pair, so $g=r$.
The component count grows with the binary width $w$ without increasing
the number of intervals.

The implementation preserves the earlier \texttt{select} method's emission
ordinal semantics.  Its additional \texttt{minimum\_intervals} property
returns the immutable canonical union as sorted disjoint half-open intervals.
\texttt{select\_minimum(rank)} selects by increasing source minimum;
\texttt{minimum\_rank(minimum)} returns that order rank and rejects points
outside the minimum set; \texttt{orbit\_rank(point)} first computes the
point's minimum and then returns its canonical rank.

\subsection{Encoded costs and the amortized improvement}

Let $E$ be the number of supplied trace events, $s$ the number of contraction
and truncation events, $g$ the total number of recorded static gaps, and $B$
a bound for integer bit lengths, including counters.  The pairing count
never increases during replay, so each contraction has at most $2k+1$ gaps
and $g=O(s(k+1))$. After independent verification, a direct list-based compiler
uses $\operatorname{poly}(E,k,B)$ bit operations, in addition to reading and
copying the explicit supplied source and certificate.  Lifting all
minimum-set hull endpoints uses at most $O(gs)$ inverse rank operations;
sorting the hulls adds $O(g\log(g+2))$ comparisons.  It retains only
$O(s+g)$ binary integers and indices, costing $O((s+g)B)$ bits.  It discards
the generator-maintenance events from the query program.

One point query uses at most $s$ forward steps and $s$ inverse rank steps.
With elementary integer arithmetic its bit cost is
\[
 O\bigl(s B^2\log(k+2)\bigr).
\]
Emission-ordinal selection uses binary search and inverse rank steps and
satisfies the same conservative bound.  Canonical minimum rank and selection
instead cost $O(B\log(g+2))$ bit work, including output arithmetic.
None of these costs depends on $N$ or $C$ as an
iteration count.  The supplied implementation is intentionally simple:
it does not claim a balanced-tree compiler or logarithmic query time in
the trace length.  On a maintained producer trace, $s$ is at most twice
the number of orbit cycles, even if the trace contains many transmission
and merger events.

For $q$ point comparisons, the complete cost is one discovery, one
independent replay and compilation, followed by
\[
 O\bigl(qs B^2\log(k+2)\bigr)
\]
bit work.  The old interface discovers two schedules per nontrivial
comparison; a stronger shared-count baseline discovers one initial schedule
and at most one augmented schedule per comparison.  The compiled interface
performs one discovery for the entire batch.  This is an exact count of
scheduled orbit computations, not a claim of a universal elapsed speedup.
If a prior count already proves the relation connected, membership is
immediate and further indexing may be unnecessary.

The documented index factory accepts only a complete independently checked
certificate, copies its input before compilation, and stores immutable
compiled data.  Later mutation of the caller's proof does not change the
index.  Invalid and unfinished traces yield no prepared index.  Discovery
retains the existing cycle allowance; verification, compilation and queries
use cooperative cancellation.  In particular, the cycle allowance is not
misrepresented as a budget for all subsequent arithmetic.

\subsection{Application to marked surfaces and signed relations}

In the maintained normal-arc model, interval orbits are precisely the
connected components of the supplied normal surface.  The compiled minimum
therefore supplies an exact, source-bound component key for any specified
edge-intersection point.  It can group an ordered list of attachment marks
by component and select a component representative without enumerating all
normal discs or all components.  These keys identify components within
one source presentation; they are not invariants under a change of
triangulation or an assertion that two attached surfaces are isomorphic.

There is also a useful signed corollary.  Suppose each pairing independently
carries a parity label, and use the existing two-sheet interval construction
on $[0,2N)$.  Let $m_2$ be a compiled minimum map for this cover.  Then the
set of attainable path parities from $x$ to $y$ is exactly
\[
 \{e\in\{0,1\}:m_2(x)=m_2(y+eN)\}.
\]
It is empty for different base components, a singleton for a component with
consistent parity, and both bits for a component containing an odd closed
walk.  Moreover, the minimum of the underlying base component is
\[
 \min\bigl(m_2(x),m_2(x+N)\bigr).
\]
Indeed, the union of the two lifted components contains both lifts of every
base point, and at least one of its two minimum labels lies in the first
sheet.  Thus one prepared cover index suffices for these point queries,
without a separate base-component search.  This corollary is a mathematical
deduction, not an implemented new signed API.  The geometric interpretation
as orientability requires a separately established orientation character;
interval reversal alone does not supply it.

This contribution makes repeated queries on one certified compressed object
cheaper.  It does not bound the number of marked states generated by a
hierarchy, compute arbitrary attachment maps, find a compressing normal
disc, or certify correspondence between a supplied triangulation and an
input knot diagram.  Those global requirements remain separate from the
local orbit-index theorem.
